\begin{longtable}{ | p{6cm} | p{7cm} | p{1.3cm} | p{1cm} | }
\hline\rowcolor{lightgray}
名称 & 描述 &  地址 & 访问 \\ \hline
\endhead
\multicolumn{4}{|l|}{\textbf{配置寄存器}} \\ \hline
\hyperref[regdesc:LPI2CANAMSTI2C0CTRLREG]{LP\_I2C\_ANA\_MST\_I2C0\_CTRL\_REG} & I2C 通路 0 控制寄存器 & 0x0000 & varies \\ \hline
\hyperref[regdesc:LPI2CANAMSTI2C0CONFREG]{LP\_I2C\_ANA\_MST\_I2C0\_CONF\_REG} & I2C 通路 0 配置寄存器 & 0x0004 & varies \\ \hline
\hyperref[regdesc:LPI2CANAMSTI2C0DATAREG]{LP\_I2C\_ANA\_MST\_I2C0\_DATA\_REG} & I2C 通路 0 数据寄存器 & 0x0008 & varies \\ \hline
\hyperref[regdesc:LPI2CANAMSTANACONF1REG]{LP\_I2C\_ANA\_MST\_ANA\_CONF1\_REG} & 模拟 I2C 输入通路使能寄存器 & 0x000C & R/W \\ \hline
\hyperref[regdesc:LPI2CANAMSTDEVICEENREG]{LP\_I2C\_ANA\_MST\_DEVICE\_EN\_REG} & 模拟 I2C 通路输出使能寄存器 & 0x0014 & R/W \\ \hline
\hyperref[regdesc:LPI2CANAMSTDATEREG]{LP\_I2C\_ANA\_MST\_DATE\_REG} & 版本控制寄存器 & 0x03FC & R/W \\ \hline
\end{longtable}
